\section{Specht modules and horizontal strips}
\label{sec:specht-tools}

We fix the tensor conventions used to identify alternating block vectors
with Specht modules. The other input needed by the induction is the
Pieri rule: it transfers constituent support across the removal of
horizontal strips. We record its dominance consequences with proofs.

All representations and tensor products are over $\C$ and are
finite-dimensional.  A partition $\lambda\vdash n$ is padded with zero
parts when necessary; its length is $\len(\lambda)$ and its conjugate is
$\lambda^t$, where $\lambda^t_j=\#\{i:\lambda_i\ge j\}$.  We identify
$\lambda$ with its diagram $\{(i,j):1\le j\le\lambda_i\}$.
For partitions of the same size, write $\lambda\unrhd\mu$ if
$\sum_{i=1}^k\lambda_i\ge\sum_{i=1}^k\mu_i$ for every $k\ge1$.
Conjugation reverses this order:
\begin{equation}\label{tools:conjugate-dominance}
 \lambda\unrhd\mu\quad\Longleftrightarrow\quad
 \mu^t\unrhd\lambda^t.
\end{equation}
Indeed, for every integer $q\ge1$,
\[
 \sum_i\max(\lambda_i-q,0)
 =\max_{k\ge0}\left(\sum_{i=1}^k\lambda_i-kq\right),
 \qquad
 \sum_{j=1}^q\lambda^t_j
 =n-\sum_i\max(\lambda_i-q,0).
\]
Dominance increases the first expression and therefore decreases the
second.  Applying the resulting implication to the conjugate partitions
proves the converse.

\subsection{Position tensors and polytabloids}

For a finite set $B$, write $S_B$ for its symmetric group. For
$\lambda\vdash |B|$, write $S_B^\lambda$ for the irreducible $S_B$-module
of type $\lambda$; omit the subscript when the position set is understood.
If $E$ has
ordered basis $e_1,\ldots,e_d$, the tensor space $E^{\otimes B}$ has word
basis $e_a=\bigotimes_{b\in B}e_{a(b)}$, indexed by maps
$a:B\to\{1,\ldots,d\}$.  The position action is
$g e_a=e_{a\circ g^{-1}}$.  All regroupings of tensor factors are ordinary
tensor identifications and introduce no signs.

Let $L$ be a set partition of $B$ whose blocks have size at most $d$.
For each block $A=(b_1,\ldots,b_r)$ choose an ordering and set
\begin{equation}\label{tools:block-tensor}
 v_A=\sum_{\pi\in S_r}\sgn(\pi)
       \bigotimes_{i=1}^r e_{\pi(i)},\qquad
 v_L=\bigotimes_{A\in L}v_A,
\end{equation}
where the $i$th factor of $v_A$ occupies position $b_i$.
Changing a block ordering changes $v_L$ only by a sign.  Its nonzero
word coefficients are $\pm1$: on each block of size $r$, the letters are
a bijection with $\{1,\ldots,r\}$.

\begin{lemma}[Alternating blocks as polytabloids]\label{lem:polytabloid}
Define $\theta_i=\#\{A\in L:|A|\ge i\}$, so that the block sizes are
the column lengths of $\theta$.  Then
$\C[S_B]v_L\cong S^\theta$.
The same assertion holds for any ordered basis of a dual vector space.
\end{lemma}

\begin{proof}
Arrange the blocks as columns of a tableau $t$ of shape $\theta$, in
nonincreasing order of their sizes, and place each ordered block from top
to bottom.  Tableau entries are the distinct elements of $B$.
A row tabloid $\{t'\}$ records the set of entries in each row.
Let $M^\theta$ be the permutation module spanned by these row tabloids.
Send $\{t'\}$ to
the word assigning letter $i$ to every entry of row $i$.
This is an $S_B$-equivariant bijection from the tabloid basis of
$M^\theta$ to the word basis of content $\theta$.
If $C_t$ is the subgroup permuting entries within each column of $t$,
the polytabloid
\[
 e_t=\sum_{c\in C_t}\sgn(c)c\{t\}
\]
maps to $v_L$: in a column, the position action replaces a permutation
of the letters by its inverse, which has the same sign.
In particular this vector is nonzero, since the original row word has
coefficient $1$.  The span of all translates of $e_t$ is the Specht
module $S^\theta$, irreducible over $\C$ by the standard polytabloid
construction \cite[Definitions~4.1 and~4.3, Theorem~4.12]{james1978}.
This proves the assertion.
The argument applies to a dual space with any chosen basis; alternatively,
an invertible basis change on that space induces an invertible map on
its position tensor power commuting with $S_B$.
\end{proof}

We use the usual complex representation-theoretic facts
\cite[Theorems~4.12 and~6.7]{james1978}: the $S^\lambda$, for
$\lambda\vdash n$, form the complete list of irreducible $S_n$-modules,
and
\begin{equation}\label{tools:dual-sign}
 (S^\lambda)^*\cong S^\lambda,
 \qquad S^\lambda\otimes\sgn\cong S^{\lambda^t}.
\end{equation}
Every complex representation of a finite group is semisimple
\cite[Proposition~1.5 and Corollary~1.6]{fultonharris1991}, so any
irreducible quotient is also a constituent.  For $H\le G$,
Frobenius reciprocity \cite[Proposition~3.17]{fultonharris1991}
takes the form
\begin{equation}\label{tools:frobenius}
 \Hom_G(\Ind_H^G A,V)
 \cong\Hom_H(A,\Res_H^G V).
\end{equation}
Induction, modeled by $\C[G]\otimes_{\C[H]}A$, is transitive and
distributes over direct sums.  The symbol
$\boxtimes$ denotes the outer tensor product for separate group actions;
ordinary tensor products of two $S_B$-modules carry the diagonal action.
The natural pairing of $E^{\otimes B}$ and $(E^*)^{\otimes B}$ satisfies
$\langle gu,gt\rangle=\langle u,t\rangle$.
Thus pairing with an invariant subspace $T$ of the dual tensor space
defines an equivariant map $u\mapsto(t\mapsto\langle u,t\rangle)$
into $T^*$.

\subsection{Pieri and the needed part of Young's rule}

For diagrams $\nu\subseteq\lambda$, the skew diagram $\lambda/\nu$
is a \emph{horizontal strip} if it has at most one box in each column.
We use the empty partition and the trivial group $S_0$ in the following
standard identity.

\begin{lemma}[Pieri rule]\label{lem:pieri}
For $\nu\vdash a$ and $b\ge0$,
\begin{equation}\label{tools:pieri}
 \Ind_{S_a\times S_b}^{S_{a+b}}
      (S^\nu\boxtimes\mathbf1)
 \cong
 \bigoplus_{\substack{\lambda\vdash a+b,\ \lambda\supseteq\nu\\
                      \lambda/\nu\text{ a horizontal strip}}}
       S^\lambda.
\end{equation}
Every summand occurs once.
\end{lemma}

This is the one-row specialization of the Littlewood--Richardson rule
\cite[Section~16, especially Rule~16.4]{james1978}: all $b$ entries of
its skew filling are $1$, so column strictness requires a horizontal
strip; the unique such filling satisfies the lattice-word condition.
We give the dominance and multiplicity
consequences needed here, including the constructive existence argument.

\begin{lemma}[Young's rule: support and diagonal multiplicity]
\label{lem:young-rule}
Let $\theta=(\theta_1,\ldots,\theta_\ell)\vdash n$. The row-tabloid
permutation module has the equivalent induced realization
\[
 M^\theta=
 \Ind_{S_{\theta_1}\times\cdots\times S_{\theta_\ell}}^{S_n}
 \mathbf1.
\]
Then $S^\tau$ occurs in $M^\theta$ if and only if
$\tau\unrhd\theta$, and $S^\theta$ occurs once.
The word module of any content obtained by reordering the parts of
$\theta$ is also isomorphic to $M^\theta$.
\end{lemma}

\begin{proof}
The case $n=0$ is immediate from the trivial representation of $S_0$.
Assume henceforth that $n>0$.
Transitivity of induction and \Cref{lem:pieri} show that the multiplicity
of $S^\tau$ counts chains
\begin{equation}\label{tools:strip-chain}
 \varnothing=\nu^{(0)}\subseteq\nu^{(1)}\subseteq\cdots
 \subseteq\nu^{(\ell)}=\tau,
 \qquad |\nu^{(k)}/\nu^{(k-1)}|=\theta_k,
\end{equation}
whose successive differences are horizontal strips.
Each strip raises the first column's height by at most one, so
$\nu^{(k)}$ has at most $k$ rows.  Since it is contained in $\tau$,
\[
 \sum_{i=1}^k\theta_i=|\nu^{(k)}|
 \le\sum_{i=1}^k\tau_i.
\]
For $k>\ell$, both initial sums equal $n$, since $\tau$ has at most
$\ell$ rows.  This proves necessity of dominance.
If $\tau=\theta$, equality forces
$\nu^{(k)}$ to be exactly the first $k$ rows of $\theta$.
These diagrams do form a chain of the required kind, proving
multiplicity one.

For sufficiency, suppose $\tau\unrhd\theta$ and induct on $\ell$.
For $\ell=1$, necessarily $\tau=(n)$.
For $\ell>1$, put $t=\theta_\ell$.
There are at least $t$ columns of $\tau$, since
$\tau_1\ge\theta_1\ge t$.  Delete a bottom box from each of the
$t$ tallest columns, choosing the rightmost columns first among equal
heights, and call the remaining diagram $\nu$.
Its column heights are still nonincreasing: unequal adjacent heights
differ by at least one, and within any group of equal heights precisely
a suffix is lowered.  Thus $\nu$ is a partition and $\tau/\nu$ is a
horizontal $t$-strip.  For every $k\ge1$ its loss in the first $k$ rows is
\begin{equation}\label{tools:prefix-loss}
 \sum_{i=1}^k\tau_i-\sum_{i=1}^k\nu_i
       =\max(0,t-\tau_{k+1}):
\end{equation}
there are $\tau_{k+1}$ columns taller than $k$, and the rule selects these
before any column whose bottom box belongs to the first $k$ rows.
If the loss is zero, the desired dominance inequality for $\nu$ follows
immediately.  If it is positive and $k<\ell$, then
\[
 \sum_{i=1}^k\nu_i
 =\sum_{i=1}^{k+1}\tau_i-t
 \ge\sum_{i=1}^{k+1}\theta_i-t
 \ge\sum_{i=1}^k\theta_i.
\]
For completeness, dominance at $\ell$ gives $\len(\tau)\le\ell$,
and dominance at $\ell-1$ gives $\tau_\ell\le t$.
All $\tau_\ell$ columns of height $\ell$ are therefore selected, so
the last row disappears and $\len(\nu)\le\ell-1$.
Hence all remaining initial sums equal the total $n-t$, and
$\nu\unrhd(\theta_1,\ldots,\theta_{\ell-1})$.
Induction gives a chain ending at $\nu$; appending the deleted strip
gives \Cref{tools:strip-chain} ending at $\tau$.

Finally, the stabilizer of a word with prescribed letter multiplicities
is the product of the symmetric groups on its equal-letter position
sets.  Reordering those multiplicities conjugates this subgroup, giving
the asserted isomorphism of permutation modules.
\end{proof}
